\SourceLeaf{385}

\par\bigskip\Needspace{8\baselineskip}
\pdfbookmark[1]{Bemerkungen zur Theorie der biquadratischen Reste}{nachlass-biquadratic-bemerkungen}
\begin{center}\large\scshape Bemerkungen.\end{center}
\BookRule{.2\textwidth}
Die Bruchstücke, die hier im Druck mit I und II bezeichnet sind, gehören nach dem Orte zu urtheilen, den die betreffenden Handschriften in einem Notizbuche einnehmen, dem Jahre $1811$ oder der zunächst folgenden Zeit an. Von den vorangehenden Versuchen, den Beweis des Fundamentaltheorems für biquadratische Reste nach den hier für den Rest $1+i$ angewandten Methoden durchzuführen, ist eine Aufzeichnung vorhanden, welche den speciellen Fall des Restes $1+2i$ erledigt und von derjenigen Bestimmung des biquadratischen Characters ausgeht, die man als Note dem Art.\ $6$ des Bruchstücks III beigefügt hat. Im übrigen lassen sich die historischen Angaben, die \textsc{Gauss} in den Anzeigen seiner arithmetischen Abhandlungen veröffentlicht hat, mit Hülfe des Nachlasses dahin ergänzen, dass die in den Artt.\ $15$ bis $20$ der \emph{Theoria residuorum biquadrat.} aufgenommenen Lehrsätze schon vor der Ausarbeitung der \emph{Theoria motus corporum coel.} niedergeschrieben sind. Die in den Anzeigen erwähnten Untersuchungen über cubische Reste werden wohl nicht zur Ausarbeitung gelangt sein; aufgezeichnet finden sich davon die mit den Hülfsmitteln, welche die Abhandlung \emph{Disquisitionum circa aequationes puras ulterior evolutio} bietet, durchgeführten Beweise der Reciprocitätssätze für zwei Primzahlen, von denen die eine reell ist.

Die Bruchstücke III bis VI bilden in der Handschrift besondere Hefte und für die drei ersten derselben weist die Form der Schriftzüge auf eine Zeit, die der für die Bruchstücke I und II nicht fern liegt, während für das letzte, Nr.\ VI, ein bedeutend späterer Zeitpunkt angenommen werden muss.

\par\medskip
[I.] Art.\ $10$. Die Bestimmung der Anzahl der Ganzenpunkte in $(z,z',z'+\zeta,z+\zeta)$ ergiebt sich aus dem Satze: bedeuten $a$ und $b$ relative Primzahlen, so geht die von $\xi+\eta i$ nach $\xi+\eta i+a+bi$ gezogene Gerade durch Einen Ganzenpunkt, wenn der imaginäre Theil von $(\xi+\eta i)\pdot(-a+bi)$ eine ganze Zahl ist.

\SourceLeaf{386}

[I.] Art.\ $17$. Die erste Umformung des letzten $S$ in dem Ausdrucke für $\Delta S$ erhält man, wenn man das betreffende Flächenstück in solche drei Theile zerlegt, dass jenes $S$ in
\[
\begin{aligned}
&S\bigl(\tfrac12m+\tfrac12-\tfrac12i,\ (\tfrac12+\tfrac12i)m,\\
&\qquad(\tfrac12+\tfrac12i)m+i,\ \tfrac12m+\tfrac12+\tfrac12i_{(+)}\bigr)\\
&-S\bigl(\tfrac12m,\ \tfrac12m+\tfrac12+\tfrac12i_{(+)},\ \tfrac12m+\tfrac12-\tfrac12i\bigr)\\
&-S\bigl((\tfrac12+\tfrac12i)m,\ \tfrac12m,\ \tfrac12m+\tfrac12-\tfrac12i\bigr)
\end{aligned}
\]
übergeht, und wenn man dann die Ganzepunkte in dem ersten Flächentheile mit Hülfe des Satzes in Art.\ $10$ auszählt und ferner berücksichtigt, dass in dem zweiten Flächentheile sich kein Ganzepunkt befindet.

Die zweite Umformung erhält man, wenn man die den Eckpunkten des dritten Flächentheils entsprechenden Grössen mit $i$ multiplicirt und um die ganze Zahl $(1-i)\frac{m-1}{2}$ vermehrt, endlich die dritte Umformung, wenn man mit der zuletzt entstandenen Figur nach Vorschrift des Art.\ $16$ diejenige vergleicht, die gegen jene die Ortsverschiedenheit $\frac{-1}{1+i}$ hat.

[I.] Art.\ $[18.]$ Eine Erläuterung zum ersten Schema findet man in dem später niedergeschriebenen hier mit [II.] bezeichneten Bruchstücke Art.\ $1$ bis $6$.

[I.] Art.\ $[18.]\ (2.)$ Die geometrische Deutung ergiebt mit Zuhülfenahme der beiden Systeme von Ganzepunkten
\[
[-2iQ-\tfrac12,\tfrac12,iQ]=[-2iQ+\tfrac12,\tfrac12,iQ]=\mathrm{IV}^{*}
\]
und
\[
-[-2iQ+\tfrac12i,-\tfrac12i,iQ]=-[-2iQ-\tfrac12i,-\tfrac12i,iQ]=\mathrm{X}^{*}
\]
die Gleichungen
\[
\begin{aligned}
\mathrm{IV}-\mathrm{IV}^{*}+[-2iQ,\tfrac12i,-\tfrac12]
&=-\mathrm{X}+\mathrm{X}^{*}+[-iQ,\tfrac12i,-\tfrac12]\\
\mathrm{XIII}+\mathrm{IV}^{*}&=[-2iQ,1,iQ]=-I(-2iQ)+I(-iQ)\\
\mathrm{V}+\mathrm{X}^{*}&=-[-2iQ,-i,iQ]=R(-2iQ)-R(-iQ)
\end{aligned}
\]
wenn allgemein $Rx$ und $Ix$ die grössten Ganzen des reellen Theils und des Coëfficienten des imaginären Theils von $x$ bedeuten.

$[-2iQ,\tfrac12i,-\tfrac12]$ ist aber die Anzahl der Ganzepunkte in dem Quadrate, dessen Mittelpunkt sich in $-2iQ-\tfrac14+\tfrac14i$ befindet und zwischen dessen Endpunkten die Ortsunterschiede $\tfrac12$ und $\tfrac12i$ Statt haben.

$[-iQ,\tfrac12i,-\tfrac12]$ oder $[-2iQ,-\tfrac12i,\tfrac12]$ ist die Anzahl der Ganzepunkte in einem gleichen Quadrate mit dem Mittelpunkte $-2iQ+\tfrac14-\tfrac14i$.

[I.] Art.\ $[18.]\ (3.)$ Mit Zuhülfenahme der Ganzepunkte $[0,\tfrac12i,-iQ]=-[-\tfrac12i,-\tfrac12i,-iQ]=\mathrm{I}^{*}$ erhält man

\SourceLeaf{387}

\[
\begin{aligned}
\mathrm{VII}-\mathrm{XII}&=\mathrm{I}-\mathrm{I}^{*}+[-iQ,\tfrac12,\tfrac12i]-[0,\tfrac12,\tfrac12i]\\
&=\mathrm{I}-\mathrm{I}^{*}+[-2iQ,\tfrac12,\tfrac12i]\\
\mathrm{II}-\mathrm{I}^{*}&=[0,-i,-iQ]=-R0+R(-iQ)\\
\mathrm{IX}-\mathrm{I}&=[0,-1,-iQ]=I0-I(-iQ)\\
\mathrm{VI}-\mathrm{XII}&=-[\tfrac12i,-1,-iQ]=-I(\tfrac12i)+I(-iQ+\tfrac12i)\\
\mathrm{VIII}-\mathrm{XIV}&=-[\tfrac12,-1-i,\tfrac12im_{(-)}]=\frac{a-1}{2}+\frac b2
\end{aligned}
\]

[II.] Art.\ $10$. Es ist
\[
\begin{aligned}
\Upsilon P'&\equiv-\Upsilon(-iP'+\tfrac12im),\\
\Upsilon P''&\equiv-1-\Upsilon(P''-\tfrac12im),\\
\Upsilon P'''&\equiv1+\Upsilon(-iP''')\ (\text{mod. }4)
\end{aligned}
\]
und $-iP'+\tfrac12im$, $P''-\tfrac12im$ sind die um $\tfrac12i$ vermehrten Ganzepunkte resp.\ in I, VI.

[III.] Art.\ $6$. Die in der Note angegebenen Regeln für die Bestimmung des $\mathrm{Dec.}\frac Mm$ habe ich der vorliegenden Abhandlung aus einem andern Orte der Handschriften beigefügt. Die erste dieser beiden Regeln, die wie leicht zu sehen mit der zweiten übereinstimmt, folgt aus der des Art.\ $6$, weil
\[
\begin{gathered}
k\equiv f\pdot i^{-n}\ (\text{mod. }m),\quad n=\Theta\frac{2fM}{m},\quad p=mm',\\
\left[\frac{2km'}p\right]^2\equiv\Theta\frac{2km'}p\ (\text{mod. }2+2i)
\end{gathered}
\]
ist.

[III.] Art.\ $8$ enthält in der Handschrift ein Beispiel zu Art.\ $7$, nemlich die Bestimmung des Decidenten von $-1+2i$ für den Modulus $-11+4i$.

[III.] Art.\ $10$. In Bezug auf die Bemerkung ‘anders auszudrücken’ kann man Art.\ $3$ des folgenden Bruchstücks [IV] vergleichen.

[IV.] Die Art.\ $1$. $2$. $4$ enthalten in der Handschrift ausser dem hier Abgedruckten noch die Anwendung auf die beiden Beispiele für $m=5+8i$, $M=9+4i$ und für $m=9+4i$, $M=5+8i$.

